%-------------------------
% Resume Builder shell: "Compact"
% Based off of: https://github.com/sb2nov/resume
% License : MIT
%
% A copy of classic.tex with ONE deliberate difference: an EXPERIENCE entry's
% heading is a SINGLE line instead of two. Classic stacks company-over-title on
% the left and dates-over-location on the right; here each pair is joined inline:
%
%   Fulton Bank, Data Scientist Intern      May 2025 – Dec 2025, Lancaster, PA
%
% That saves a line per experience entry while keeping the classic look otherwise
% identical. See \resumeSubheading below.
%
% EDUCATION IS DELIBERATELY EXCLUDED. Degree titles are long — "Bachelor's of
% Technology in Computer Science and Engineering (Spl. in AI and ML)" — so
% collapsing them onto one line forces an ugly mid-phrase wrap that orphans a
% fragment under the school name. Education therefore keeps classic's two-line
% layout via \resumeEducationSubheading, which the renderer emits for education
% entries specifically (other shells alias it to \resumeSubheading).
%
% This file is the rendering shell for the Resume Builder. Everything in the
% preamble (packages, margins, \resume* custom commands) is fixed; only the
% body sections between \begin{document} and \end{document} are produced from
% structured data by backend/resume_render.py. Do NOT hand-edit the {{...}}
% placeholders — they are filled in programmatically.
%------------------------

\documentclass[letterpaper,10pt]{article}

\usepackage{latexsym}
\usepackage[empty]{fullpage}
\usepackage{titlesec}
\usepackage{marvosym}
\usepackage[usenames,dvipsnames]{color}
\usepackage{verbatim}
\usepackage{enumitem}
\usepackage[hidelinks]{hyperref}
\usepackage{fancyhdr}
\usepackage[english]{babel}
\usepackage{tabularx}
% array provides the >{...} column modifier and \arraybackslash used by the
% single-line \resumeSubheading below; tabularx alone does not define them.
\usepackage{array}
% NOTE: fontawesome5 was in the source template but no \fa… icons are used.
% It requires loadable OTF fonts that Tectonic's XeTeX can't resolve in this
% environment (aborts with SIGABRT), so it's intentionally omitted.
\usepackage{multicol}
\usepackage{setspace}
\setstretch{1.02}
\setlength{\multicolsep}{-3.0pt}
\setlength{\columnsep}{-1pt}
% glyphtounicode / \pdfgentounicode are pdfTeX-only primitives used to make the
% PDF ATS-parsable. Tectonic compiles with XeTeX (which emits unicode-mapped
% PDFs natively and does NOT define these), so we load them only when running
% under pdfTeX. Without this guard XeTeX aborts on an undefined control sequence.
\ifdefined\pdfgentounicode
  \input{glyphtounicode}
\fi

\pagestyle{fancy}
\fancyhf{}
\fancyfoot{}
\renewcommand{\headrulewidth}{0pt}
\renewcommand{\footrulewidth}{0pt}

% Optimized margins for professional appearance
\addtolength{\oddsidemargin}{-0.55in}
\addtolength{\evensidemargin}{-0.55in}
\addtolength{\textwidth}{1.1in}
\addtolength{\topmargin}{-0.65in}
\addtolength{\textheight}{1.4in}

\urlstyle{same}

\raggedbottom
\raggedright
\setlength{\tabcolsep}{0in}

% Section heading look — the original: a \titlerule with a trailing \vspace{-4pt}
% that sets the rule weight and the gap below the heading. The rule thickness is
% stated EXPLICITLY (0.4pt) rather than left to \titlerule's default, which
% inherits \arrayrulewidth and can render heavier than intended. The rule MUST
% be brace-wrapped — titlesec delimits its [after] option with brackets, so a
% bare \titlerule[0.4pt] inside it ends the option early and the compile dies
% with "Missing \begin{document}". `before`=0 because the
% gap ABOVE a section is owned by \sectionsep (below), not by titlespacing;
% `after`=8pt with the -4pt nets the original below-rule spacing.
\titleformat{\section}{
  \scshape\raggedright\large\bfseries
}{}{0em}{}[{\color{black}\titlerule[0.4pt]} \vspace{-4pt}]
\titlespacing*{\section}{0pt}{0pt}{8pt}

% Uniform inter-section spacing (order-independent). The gap ABOVE every section
% is emitted by \sectionsep — placed BEFORE each section block (including the
% first) by the renderer — rather than baked into each section relative to its
% neighbour. \addvspace MERGES with any pending vertical glue instead of stacking,
% so residual glue can't double the gap. This works together with two things in
% resume_render.py: (a) sections don't emit a trailing negative \vspace on their
% LAST entry (which used to leak into the next section and is why gaps were
% uneven), and (b) the header's -8pt is offset just below so the first section
% matches the rest. Tune \sectiongap to retune all inter-section gaps at once.
\newlength{\sectiongap}\setlength{\sectiongap}{7pt}
\newcommand{\sectionsep}{\par\addvspace{\sectiongap}}

% Ensure that generate pdf is machine readable/ATS parsable (pdfTeX only)
\ifdefined\pdfgentounicode \pdfgentounicode=1 \fi

%-------------------------
% Custom commands with balanced spacing
\newcommand{\resumeItem}[1]{
  \item\small{
    {#1 \vspace{-2pt}}
  }
}

% Single-line entry heading — the ONE thing that differs from classic.tex.
% Args are (#1 company/school, #2 dates, #3 title/degree, #4 location), same as
% classic. Rather than stacking #1 over #3 and #2 over #4 on two rows, each pair
% is joined on one row with a comma, keeping classic's weighting: organisation
% and dates bold, role and location italic.
%
% \resumePair joins two fields as "a, b" — but emits just "a" when b is blank and
% just "b" when a is blank, so neither a missing role nor a missing location can
% leave a dangling or leading comma. Both cases are real: the renderer always
% passes all four fields, empty string included, and an entry may legitimately
% have no title (#1) or no location (#4).
%
% `\ifx\relax#1\relax` is the standard empty-argument test: it expands to
% \ifx\relax\relax (true) only when #1 is blank.
\newcommand{\resumePair}[2]{%
  \ifx\relax#1\relax
    #2%
  \else
    \textbf{#1}\ifx\relax#2\relax\else\textit{\small, #2}\fi
  \fi
}
%
% WHY tabularx AND NOT tabular*: classic's `tabular*` with \extracolsep{\fill}
% distributes only LEFTOVER space, and neither cell can wrap. On one line the
% left cell (organisation + full degree title) is often far wider than half the
% page — e.g. "Master of Science in Data Science (CCAS Dean's Award)" — so the
% fill collapses to zero, the right cell butts straight against the left text,
% and the row runs off the page edge. tabularx instead gives the right column its
% NATURAL width and lets the X (left) column absorb what remains, wrapping onto a
% second line when the content genuinely doesn't fit. A too-long entry therefore
% degrades to two lines (still no worse than classic) rather than overflowing.
%
% \raggedright in the X column suppresses justification (which would otherwise
% stretch interword space across a wrapped line); the \arraybackslash restores \\
% after \raggedright redefines it. The right column is `r` so dates stay flush to
% the margin exactly as in classic.
\newcommand{\resumeSubheading}[4]{
  \vspace{-2pt}\item
    \begin{tabularx}{1.0\textwidth}[t]{@{}>{\raggedright\arraybackslash}X@{\hspace{1em}}r@{}}
      \resumePair{#1}{#3} &
      \textbf{\small #2}\ifx\relax#4\relax\else\textit{\small, #4}\fi \\
    \end{tabularx}\vspace{-7pt}
}

\newcommand{\resumeProjectHeading}[2]{
    \item
    \begin{tabular*}{1.001\textwidth}{l@{\extracolsep{\fill}}r}
      \small#1 & \textbf{\small #2}\\
    \end{tabular*}\vspace{-7pt}
}

% Education keeps CLASSIC's two-line layout — this is classic.tex's
% \resumeSubheading verbatim. See the note at the top of this file for why
% education is excluded from the single-line treatment.
\newcommand{\resumeEducationSubheading}[4]{
  \vspace{-2pt}\item
    \begin{tabular*}{1.0\textwidth}[t]{l@{\extracolsep{\fill}}r}
      \textbf{#1} & \textbf{\small #2} \\
      \ifx\relax#3\relax\ifx\relax#4\relax\else
      \textit{\small#3} & \textit{\small #4} \\\fi\else
      \textit{\small#3} & \textit{\small #4} \\\fi
    \end{tabular*}\vspace{-7pt}
}

\renewcommand\labelitemi{$\vcenter{\hbox{\tiny$\bullet$}}$}
\renewcommand\labelitemii{$\vcenter{\hbox{\tiny$\bullet$}}$}

\newcommand{\resumeSubHeadingListStart}{\begin{itemize}[leftmargin=0.0in, label={}]}
\newcommand{\resumeSubHeadingListEnd}{\end{itemize}}
\newcommand{\resumeItemListStart}{\begin{itemize}}
\newcommand{\resumeItemListEnd}{\end{itemize}\vspace{-5pt}}

%-------------------------------------------
%%%%%%  RESUME STARTS HERE  %%%%%%%%%%%%%%%%%%%%%%%%%%%%

\begin{document}

%----------HEADING----------
\begin{center}
    {\Huge \scshape {{FULL_NAME}}} \\ \vspace{3pt}
    \small {{CONTACT_LINE}}
    \vspace{-8pt}
\end{center}
% Partially offset the header's trailing -8pt so the gap above the FIRST section
% matches the others (measured: this lands the first-section gap on the same
% baseline as the \sectionsep gap the body sections get).
\par\vspace{1pt}

{{SECTIONS}}

\end{document}
